% helden/dorle.tex — Werte eines Helden. Nur Werte, keine Koordinaten.
% Fehlt ein Wert, bleibt das Feld leer. Kein Abbruch.
% Die Feldnamen sind die der Feldtabelle und in beiden Quellfassungen gleich.

\wert{s1_name}{Dorle Griesgram}
\wert{s1_geschlecht}{weiblich}
\wert{s1_kultur}{Mittelländisches Dorf}
% s1_profession absichtlich nicht gesetzt — Probe auf den Leerfall

% Abgeleitete Werte. Der Bogen hat die Formeln aufgedruckt; hier stehen nur
% Zahlen, es wird nichts nachgerechnet.
\wert{s1_lep_wert}{32}
\wert{s1_lep_bonus}{2}
\wert{s1_lep_max}{34}
\wert{s1_sk_wert}{1}
\wert{s1_sk_max}{1}
\wert{s1_zk_wert}{2}
\wert{s1_zk_max}{2}
\wert{s1_aw_wert}{5}
\wert{s1_aw_max}{5}
\wert{s1_schips_wert}{3}
\wert{s1_schips_max}{3}
% AsP und KaP absichtlich leer — Dorle ist weder Zauberin noch Geweihte

% Restliche Angaben Seite 1
\wert{s1_spezies}{Mensch}
\wert{s1_sozialstatus}{unfrei}
\wert{s1_heimatort}{Fuchsgrund}
\wert{s1_alter}{34}
\wert{s1_haarfarbe}{graubraun}
\wert{s1_augenfarbe}{grau}
\wert{s1_groesse}{1,64 m}
\wert{s1_gewicht}{58 kg}
% s1_familie, s1_charakter_* absichtlich leer

% --- Bogen 3: Kampfwerte ---------------------------------------------------
\wert{s3_gs}{8}
\wert{s3_le}{32}
\wert{s3_aw}{5}
\wert{s3_ini}{12}
\wert{s3_sk}{1}
\wert{s3_zk}{2}

\wert{s3_kt_dolche_ktw}{10}
\wert{s3_kt_dolche_atfk}{8}
\wert{s3_kt_dolche_pa}{5}
\wert{s3_kt_raufen_ktw}{8}
\wert{s3_kt_raufen_atfk}{7}
\wert{s3_kt_raufen_pa}{4}
\wert{s3_kt_armbrueste_ktw}{6}
\wert{s3_kt_armbrueste_atfk}{6}
\wert{s3_kt_schwerter_ktw}{12}
\wert{s3_kt_schwerter_atfk}{9}
\wert{s3_kt_schwerter_pa}{6}
\wert{s3_kt_wurfwaffen_ktw}{7}
\wert{s3_kt_wurfwaffen_atfk}{7}

\wert{s3_ksf_1}{Wuchtschlag I, Finte I}
\wert{s3_ksf_2}{Klingensturm}

\wert{s3_lep_max}{34}
\wert{s3_lep_aktuell}{34}
\wert{s3_lep_schwelle_1}{25}
\wert{s3_lep_schwelle_2}{17}
\wert{s3_lep_schwelle_3}{8}
\wert{s3_lep_schwelle_4}{5}

% --- Bogen 3: Tabellen (Testwerte zur Sichtprobe der Spalten) -------------
\wert{s3_nk_1_waffe}{Dolch}
\wert{s3_nk_1_technik}{Dolche}
\wert{s3_nk_1_schaden}{KK 14/3}
\wert{s3_nk_1_tp_basis}{4}
\wert{s3_nk_1_tp_gesamt}{5}
\wert{s3_nk_1_mod}{0/0}
\wert{s3_nk_1_reichweite}{kurz}
\wert{s3_nk_1_at}{9}
\wert{s3_nk_1_pa}{5}
\wert{s3_nk_1_gewicht}{1 Unz}

\wert{s3_fk_1_waffe}{Schwere Armbrust}
\wert{s3_fk_1_technik}{Armbrueste}
\wert{s3_fk_1_ladezeit}{16 Akt}
\wert{s3_fk_1_tp}{13}
\wert{s3_fk_1_reichweite}{25/50/100}
\wert{s3_fk_1_fernkampf}{6}
\wert{s3_fk_1_munition}{12 Bolzen}
\wert{s3_fk_1_gewicht}{5 Stn}

\wert{s3_ru_1_name}{Lederharnisch}
\wert{s3_ru_1_rs}{3}
\wert{s3_ru_1_be}{1}
\wert{s3_ru_1_abzuege}{-}
\wert{s3_ru_1_gewicht}{3 Stn}
\wert{s3_ru_1_wo}{Torso}

\wert{s3_sch_1_name}{Holzschild}
\wert{s3_sch_1_struktur}{4}
\wert{s3_sch_1_mod}{0/+1}
\wert{s3_sch_1_gewicht}{2 Stn}

\wert{s3_zu_belastung_stufe1}{X}
\wert{s3_zu_belastung_stufe2}{X}
\wert{s3_zu_belastung_stufe3}{X}
\wert{s3_zu_belastung_stufe4}{X}
\wert{s3_zu_frei_stufe1}{X}

% --- Bogen 1: Eigenschaften -----------------------------------------------
\wert{s1_mu}{12}
\wert{s1_kl}{11}
\wert{s1_in}{13}
\wert{s1_ch}{10}
\wert{s1_ff}{12}
\wert{s1_ge}{13}
\wert{s1_ko}{12}
\wert{s1_kk}{11}

% --- Bogen 2: Talente (Stichproben in allen fünf Gruppen) ----------------
\wert{s2_belastung}{1}
\wert{s2_mu}{12} \wert{s2_kl}{11} \wert{s2_in}{13} \wert{s2_ch}{10}
\wert{s2_ff}{12} \wert{s2_ge}{13} \wert{s2_ko}{12} \wert{s2_kk}{11}
\wert{s2_fliegen_fw}{0}
\wert{s2_klettern_fw}{6}
\wert{s2_klettern_rpr}{-}
\wert{s2_klettern_anm}{Fels}
\wert{s2_wildnisleben_fw}{9}
\wert{s2_wildnisleben_rpr}{3}
\wert{s2_wildnisleben_anm}{Wald}
\wert{s2_willenskraft_fw}{4}
\wert{s2_brettspiel_fw}{3}
\wert{s2_sternkunde_fw}{2}
\wert{s2_alchimie_fw}{1}
\wert{s2_stoffbearbeitung_fw}{7}
\wert{s2_stoffbearbeitung_rpr}{2}
\wert{s2_stoffbearbeitung_anm}{Flicken}
\wert{s2_mod_mu_0}{X}
\wert{s2_mod_kk_p2}{X}
\wert{s2_sprachen_1}{Garethi (Muttersprache)}
\wert{s2_sprachen_2}{Zwergisch (Grundkenntnisse)}
\wert{s2_schriften_1}{Kusliker Zeichen}

% --- Bogen 4: Liturgien ---------------------------------------------------
\wert{s4_kap_max}{0}
\wert{s4_kap_aktuell}{0}
\wert{s4_lit_1_name}{Bannstrahl}
\wert{s4_lit_1_probe}{MU/KL/CH}
\wert{s4_lit_1_fw}{6}
\wert{s4_lit_1_kosten}{4 KaP}
\wert{s4_lit_1_dauer}{2 Akt}
\wert{s4_lit_1_reichweite}{16 Schritt}
\wert{s4_lit_1_wirkdauer}{QS Runden}
\wert{s4_lit_1_aspekt}{Bann}
\wert{s4_lit_1_stf}{B}
\wert{s4_lit_1_wirkung}{Ziel erhaelt -1 auf Proben}
\wert{s4_lit_1_seite}{9}
\wert{s4_lit_20_name}{letzte Zeile}
\wert{s4_tradition}{Kirche des Praios}
\wert{s4_leiteigenschaft}{CH}
\wert{s4_aspekt}{Handel}
\wert{s4_talente_1}{Etikette, Rechtskunde}
\wert{s4_zeremonialgegenstand_1}{Sonnenscheibe}
\wert{s4_ksf_1}{Liturgiestil}
\wert{s4_segnungen_1}{Segen des Lichts}

% --- Bogen 5: Zauber ------------------------------------------------------
\wert{s5_asp_max}{0}
\wert{s5_zau_1_name}{Flim Flam}
\wert{s5_zau_1_probe}{KL/CH/FF}
\wert{s5_zau_1_fw}{8}
\wert{s5_zau_1_kosten}{2 AsP}
\wert{s5_zau_1_dauer}{2 Akt}
\wert{s5_zau_1_reichweite}{16 Schritt}
\wert{s5_zau_1_wirkdauer}{QS Minuten}
\wert{s5_zau_1_merkmal}{Illusion}
\wert{s5_zau_1_stf}{A}
\wert{s5_zau_1_wirkung}{erzeugt Lichterscheinung}
\wert{s5_zau_1_seite}{271}
\wert{s5_zau_20_name}{letzte Zeile}
\wert{s5_tradition}{Gildenmagie}
\wert{s5_merkmal}{Illusion}
\wert{s5_artefakt_1}{Stab}
\wert{s5_msf_1}{Merkmalskenntnis}
\wert{s5_tricks_1}{Farbenspiel}

% --- Bogen 6: Ausrüstung --------------------------------------------------
\wert{s6_gg_l_1_gegenstand}{Wanderstab}
\wert{s6_gg_l_1_gewicht}{2 Stn}
\wert{s6_gg_l_1_wo}{Hand}
\wert{s6_gg_l_30_gegenstand}{letzte Zeile links}
\wert{s6_gg_r_1_gegenstand}{Wollmantel}
\wert{s6_gg_r_1_gewicht}{3 Stn}
\wert{s6_gg_r_1_wo}{getragen}
\wert{s6_gg_r_30_gegenstand}{letzte Zeile rechts}
\wert{s6_gesamtgewicht_l}{14 Stn}
\wert{s6_gesamtgewicht_r}{9 Stn}
\wert{s6_tragkraft_l}{22}
\wert{s6_tragkraft_r}{22}
\wert{s6_tk_basis}{22}
\wert{s6_tk_4}{26}
\wert{s6_tk_8}{30}
\wert{s6_tk_12}{34}
\wert{s6_tk_16}{38}
\wert{s6_dukaten}{2}
\wert{s6_silbertaler}{14}
\wert{s6_heller}{7}
\wert{s6_kreuzer}{3}
\wert{s6_tier_name}{Grauschnauze}
\wert{s6_tier_lo}{2}
\wert{s6_tier_ap}{80}
\wert{s6_tier_typus}{Haushund}
\wert{s6_tier_lep}{18}
\wert{s6_tier_mu}{13}
\wert{s6_tier_kl}{6}
\wert{s6_tier_in}{14}
\wert{s6_tier_ch}{11}
\wert{s6_tier_ff}{8}
\wert{s6_tier_ge}{16}
\wert{s6_tier_ko}{13}
\wert{s6_tier_kk}{11}
\wert{s6_tier_sk}{2}
\wert{s6_tier_zk}{3}
\wert{s6_tier_rs}{0}
\wert{s6_tier_ini}{15}
\wert{s6_tier_gs}{12}
\wert{s6_tier_angriff}{Biss}
\wert{s6_tier_at}{11}
\wert{s6_tier_vw}{5}
\wert{s6_tier_tp}{1W6}
\wert{s6_tier_rw}{kurz}
\wert{s6_tier_aktionen}{1 Angriff}
\wert{s6_tier_sf}{Gutes Gehoer}

% Eigenschaftsleiste auf Bogen 4 und 5 (zur Prüfung der Zellposition)
\wert{s4_mu}{12} \wert{s4_kl}{11} \wert{s4_in}{13} \wert{s4_ch}{10}
\wert{s4_ff}{12} \wert{s4_ge}{13} \wert{s4_ko}{12} \wert{s4_kk}{11}
\wert{s5_mu}{12} \wert{s5_kl}{11} \wert{s5_in}{13} \wert{s5_ch}{10}
\wert{s5_ff}{12} \wert{s5_ge}{13} \wert{s5_ko}{12} \wert{s5_kk}{11}

% Bogen 1: Textblöcke
\wert{s1_vorteile_1}{Zäher Hund, Ortskenntnis (Fuchsgrund)}
\wert{s1_nachteile_1}{Streitsucht, Neugier}
\wert{s1_sf_1}{Ortskenntnis, Wildnisleben-Anwendung}
\wert{s1_schips_aktuell}{3}
\wert{s1_erfahrungsgrad}{Erfahren}
\wert{s1_ap_gesamt}{1100}
\wert{s1_ap_gesammelt}{1100}
\wert{s1_ap_ausgegeben}{1050}

% Heldenbild. Pfad relativ zum Projektverzeichnis. Fehlt die Datei, bleibt
% der Rahmen leer und der Lauf gibt nur einen Hinweis ins Log.
\wert{s1_bild}{helden/bilder/pruefbild.png}

% --- Charaktermappe --------------------------------------------------------
% Der Umschlag (mappe/mappe.tex) zieht seine Angaben aus derselben Datei.
% Fehlt eine davon, laesst die Mappe die Stelle leer -- bis auf den Namen,
% den sie notfalls von s1_name nimmt.

% Akzentfarbe: sechs Hexstellen, ohne Raute. Sie faerbt Rahmen, Name und
% Profession. Am Vorbild gemessen: Hexe AD6ED0, Streuner D88667.
\wert{mappe_akzentfarbe}{6F9455}

% Name und Profession. Ohne mappe_name nimmt die Mappe s1_name -- hier
% ueberschrieben, weil der Titel zweizeilig gesetzt werden soll.
\wert{mappe_name}{DORLE\\GRIESGRAM}
\wert{mappe_profession}{KRAEUTERSAMMLERIN}

% Ganzkoerperbild fuer Titel und Wasserzeichen der Rueckseite. Ohne
% mappe_bild nimmt die Mappe s1_bild.
\wert{mappe_bild}{helden/bilder/pruefbild.png}

% Kurzempfehlung im Fusskasten der Titelseite.
\wert{mappe_empfehlung}{Diese Heldin solltest du spielen, wenn du eine
  bodenstaendige Kraeutersammlerin sein willst, die sich in Wald und Feld
  auskennt, jedem Dorf ein Mittel gegen das Fieber mischen kann und lieber
  einmal zuviel nachfragt als einmal zuwenig.}

% Rueckseite: Hintergrund. Leerzeilen trennen Absaetze.
\wert{mappe_hintergrund}{Dorle Griesgram wurde in Fuchsgrund geboren und ist
  dort nie lange fortgeblieben. Ihre Mutter hat ihr gezeigt, welche Wurzel
  gegen Husten hilft und welche das Vieh umbringt, und beides hat sie
  behalten.

  Sie ist keine Kaempferin. Was sie kann, kann sie gruendlich: Kraeuter
  finden, Wunden versorgen, mit Tieren umgehen und den Weg durch den Wald
  auch dann noch wissen, wenn der Nebel steht. Wer mit ihr reist, kommt
  seltener krank an und oefter ueberhaupt.

  Ihren Beinamen traegt sie zu Recht. Dorle redet wenig, und wenn sie redet,
  dann selten Erfreuliches. Wer sie eine Weile kennt, merkt aber, dass sie
  murrt, waehrend sie hilft -- und dass sie nie aufhoert zu helfen.}

% Rueckseite: freier Kasten am Fuss.
\wert{mappe_fusstitel}{ICH WUERDE GERNE EINEN KRAEUTERSAMMLER SPIELEN}
\wert{mappe_fusstext}{Wenn du zwar gerne eine bodenstaendige
  Kraeuterkundige spielen moechtest, diese aber ein Mann sein soll, dann ist
  das kein Problem. Aendere einfach den Namen zu Dorlan Griesgram, und
  fertig ist dein neuer Kraeutersammler.}

% Linkes Innenblatt: unter dem Pflichttext ist Platz fuer einen eigenen
% Kasten. Ohne diese Werte bleibt die Flaeche leer.
\wert{mappe_innen_links_titel}{ZUM ABENTEUER FUCHSGRUND}
\wert{mappe_innen_links_text}{Dorle ist im Fuchsgrund zu Hause. Wer sie
  spielt, kennt die Leute im Dorf beim Namen und weiss, wo der Bach im
  Sommer versiegt -- das ist ein Vorsprung und eine Verpflichtung
  zugleich.}

% Rechtes Innenblatt: ohne mappe_innen_rechts_text erscheinen Notizlinien.
\wert{mappe_innen_rechts_titel}{NOTIZEN UND BEKANNTSCHAFTEN}

% Copyright-Zeile im Pflichttext.
\wert{mappe_copyright}{2026 Leif Janzik}
